% PAGES: 303-305
% First chunk of Chapter 10 (the book's final chapter). Nothing open from
% Chapter 9 -- this chunk starts a fresh chapter, and the orchestrator writes
% ch10.tex with the \chapter{} and \label{} commands itself (matching the
% ch01-ch06/ch08 pattern), so this file does not repeat them.

\noindent Determinants, permutation matrices, orthogonal vectors and orthogonal matrices,
quadratic forms.
\medskip

\noindent Multivariable functions.
\medskip

\noindent Lagrange multipliers.
\medskip

\noindent The ``Big O'' notation.
\medskip

\noindent European options overview.
\medskip

\noindent Eigenvalues of symmetric matrices.
\medskip

\noindent Row rank equal to column rank.
\medskip

\noindent Technical results for Cholesky and LU decompositions.
\medskip

\section{Numerical linear algebra tools.}

In this section, we give a brief overview to determinants, permutation matrices,
orthogonal vectors and orthogonal matrices, and quadratic forms.

\subsection{Determinants}

The concept of the determinant of a square matrix is of little relevance for numerical
linear algebra purposes.\footnote{The condition number of a matrix is much more relevant,
since, it indicates not only whether a matrix is nonsingular, as does the determinant of
the matrix, but also how far from being nonsingular (i.e., ill--conditioned or
well--conditioned) a matrix is; see Demmel \cite{13} for more details.} However, if
needed, the determinant of the matrix can be computed efficiently using the LU
decomposition of the matrix, which requires $\frac23 n^3 + O(n^2)$ operations. Using the
classical definition of the determinant to compute it would require a huge number of
operations, on the order of $n \cdot n! \approx \frac{n^{n+1}}{e^n}$.

A very important use of the concept of the determinant of a matrix is as an equivalent
characterization of a matrix being nonsingular if and only if the determinant of the
matrix is nonzero; see Theorem 1.2.

The determinant of a diagonal matrix $D$ is equal to the product of all the main
diagonal entries of the matrix, i.e.,
\begin{equation}
\det(D) \;=\; \prod_{i=1}^n D(i,i).
\end{equation}

In particular, the determinant of the identity matrix is 1, i.e.,
\begin{equation}
\det(I) \;=\; 1.
\end{equation}

The determinant of a lower triangular matrix $L$ is equal to the product of all the
main diagonal entries of the matrix, i.e.,
\begin{equation}
\det(L) \;=\; \prod_{i=1}^n L(i,i).
\end{equation}

The determinant of an upper triangular matrix $U$ is equal to the product of all the
main diagonal entries of the matrix, i.e.,
\begin{equation}
\det(U) \;=\; \prod_{i=1}^n U(i,i).
\end{equation}

Since the transpose of an upper triangular matrix is lower triangular, it follows from
(10.3) and (10.4)
\begin{equation}
\det(U^t) \;=\; \prod_{i=1}^n U(i,i) \;=\; \det(U).
\end{equation}

The determinant of the $2\times2$ matrix $A = \begin{pmatrix} a & b \\ c & d
\end{pmatrix}$ is
\begin{equation}
\det(A) \;=\; ad \;-\; bc.
\end{equation}

Note that the inverse of the $2\times2$ matrix $A$ is
\begin{equation}
A^{-1} \;=\; \frac{1}{\det(A)}\begin{pmatrix} d & -b \\ -c & a \end{pmatrix}
\;=\; \frac{1}{ad-bc}\begin{pmatrix} d & -b \\ -c & a \end{pmatrix}.
\end{equation}

The determinant of the $3\times3$ matrix $A = \begin{pmatrix} a_{1,1} & a_{1,2} &
a_{1,3} \\ a_{2,1} & a_{2,2} & a_{2,3} \\ a_{3,1} & a_{3,2} & a_{3,3} \end{pmatrix}$ is
\begin{align}
\det(A) \;&=\; a_{1,1}a_{2,2}a_{3,3} \;+\; a_{1,2}a_{2,3}a_{3,1} \;+\; a_{1,3}a_{2,1}a_{3,2} \notag \\
&\quad -\; a_{1,3}a_{2,2}a_{3,1} \;-\; a_{1,1}a_{2,3}a_{3,2} \;-\; a_{1,2}a_{2,1}a_{3,3}.
\end{align}

\begin{lemma}
Let $A$ and $B$ be square matrices of the same size. Then,
\begin{equation}
\det(AB) \;=\; \det(A)\,\det(B).
\end{equation}
\end{lemma}

\begin{lemma}
Let $A$ be a square matrix. Then,
\[
\det(A^t) \;=\; \det(A).
\]
\end{lemma}

\subsection{Permutation matrices}

\begin{definition}
A permutation matrix is a matrix obtained by permuting the columns (or, equivalently,
the rows) of the identity matrix.
\end{definition}

For every permutation matrix $P$ of size $n$ there exists a permutation\footnote{A
permutation of the numbers $1,2,\ldots,n$ is a function $\tau : \{1,2,\ldots,n\} \to
\{1,2,\ldots,n\}$ which is one--to--one and onto, i.e., a function that assigns to every
number $i$ between $1$ and $n$ a unique number $\tau(i)$ between $1$ and $n$ such that,
if $1 \le i \ne j \le n$, then $\tau(i) \ne \tau(j)$.} $\tau_c$ of the numbers $1,2,
\ldots,n$ such that
\[
P \;=\; \mathrm{col}\,\!\left(e_{\tau_c(k)}\right)_{k=1:n}.
\]

Similarly, for every permutation matrix $P$ of size $n$ there exists a permutation
function $\tau_r$ of the numbers $1:n$ such that
\[
P \;=\; \mathrm{row}\,\!\left(e_{\tau_r(j)}^t\right)_{j=1:n}.
\]

However, the column and row permutation functions corresponding to the same
permutation matrix need not be the same, as seen in the example below:

\begin{examplex}
The matrix
\[
P \;=\; \begin{pmatrix}
0 & 0 & 1 & 0 & 0 \\
1 & 0 & 0 & 0 & 0 \\
0 & 0 & 0 & 0 & 1 \\
0 & 1 & 0 & 0 & 0 \\
0 & 0 & 0 & 1 & 0
\end{pmatrix}
\]
is a permutation matrix.

The row form of $P$ is
\[
P \;=\; \begin{pmatrix} e_3^t \\ e_1^t \\ e_5^t \\ e_2^t \\ e_4^t \end{pmatrix}
\;=\; \begin{pmatrix} 3 \\ 1 \\ 5 \\ 2 \\ 4 \end{pmatrix},
\]
and can be expressed as
\[
P \;=\; \mathrm{row}\,\!\left(e_{\tau_r(j)}^t\right)_{j=1:5},
\]
where $\tau_r : \{1,2,3,4,5\} \to \{1,2,3,4,5\}$ is given by
\[
\tau_r(1) = 3; \quad \tau_r(2) = 1; \quad \tau_r(3) = 5; \quad \tau_r(4) = 2; \quad
\tau_r(5) = 4.
\]

The column form of $P$ is
\[
P \;=\; (e_2 \mid e_4 \mid e_1 \mid e_5 \mid e_3) \;=\; (2\ 4\ 1\ 5\ 3),
\]
and can be expressed as
\[
P \;=\; \mathrm{col}\,\!\left(e_{\tau_c(k)}\right)_{k=1:5},
\]
where $\tau_c : \{1,2,3,4,5\} \to \{1,2,3,4,5\}$ is given by
\[
\tau_c(1) = 2; \quad \tau_c(2) = 4; \quad \tau_c(3) = 1; \quad \tau_c(4) = 5; \quad
\tau_c(5) = 3.
\]
\end{examplex}
